\scanpage{134}
\mspage{V.6}
% archive sheet number: 133
% manuscript lines 1-3
2. \uline{Application du th. de Lefschetz affine.}

2.1. Soit \deleted{\unsure{$a$}} $x\in X_0$, et
\[
n_x=\deleted{\dim X_{\unsure{\eta|y}}\unsure{\text{ local}}}
\sup_{\substack{X_i\text{ comp. irréduct.}\\\text{de }X\text{ passant par }x}}\dim X_{i\eta}
\quad\deleted{\unsure{i\in I}}
\tag*{(2.1.0)}
\]
\inserted{Supposons $U\to X_\eta$ affine.}
% manuscript lines 4-7
Alors dans (1.4.2) les $X_{\alpha\bar\eta}$ sont des affines et \unsure{finis|quasi-}\deleted{\unsure{compacts $U_{\bar\eta}$}} étales sur $X_\eta$, \inserted{et si on veut au dessus de $\bigcup_i X_{i\eta}$;} donc affines de dim $\leq n$, donc leur $H^i$ est nul pour $i>n$ \inserted{(SGA \deleted{\unsure{4}}4 \uline{XIV} ---)}. On trouve donc
\[
(\Psi^i_{\unsure{\nu|v|*}})_{\bar x}=0\qquad\text{pour }i>n_x.
\tag*{(2.1.\unsure{1})}
\]
% manuscript lines 8-11
On va améliorer ce résultat en prouvant que
\[
(\Psi^i_\nu)_{\bar x}=0\qquad\text{pour }i>n_x-\dim\overline{\{x\}}.
\tag*{(2.1.2)}
\]
ce qui implique aussi
\[
(\Psi^i_\nu)_{\bar x}=0\qquad\text{si }i>\dim\underline{\unsure{\mathcal O|\mathscr O}}_{X_0,x}=\dim_x(X_0)-\dim\overline{\{x\}}
\tag*{(2.1.3)}
\]
% manuscript lines 12-18
2.2. Pour ceci, notons qu'on peut \deleted{\unsure{supposer (sans changer}} remplacer $X$ par \deleted{$X_{\unsure{\bar\eta}}$} l'adhérence schématique $\overline{(X_\eta)_{\unsure{\mathrm{réd}|\mathrm{red}}}}$ [sans changer les \deleted{\unsure{valeurs des}} $\Psi^{(i)}_\nu$ \deleted{\unsure{\}}} ni $n_x$, et en diminuant $\dim_x X_0$], alors \deleted{\unsure{les deux} $n_x-\dim_x X_0$, \unsure{et les deux}} \deleted{\unsure{de}} $n_x=\dim_x X_0$ \deleted{$\dim X_0\leq\unsure{n}$} et les deux formules (2.1.1) et (2.1.2) sont identiques.

\deleted{\unsure{quitte à remplacer $X$ par un voisinage ouvert affine de $x$ dans $X$,}}

% manuscript lines 19-25: entire block cancelled by long diagonal strokes
\deleted{Alors, il existe un morphisme fini}
\[
\deleted{X\xrightarrow{h}Y=S[T_1,\ldots,T_n],}
\]
\deleted{et si on pose}
\[
\deleted{F=h_*(\Lambda_X),\quad y=h(x),}
\]
\deleted{on trouve que les faisceaux $\Psi^{(i)}_{\unsure{f|g}}(F)$ relatifs à $F$ (qui se définissent de façon évidente, comme}

% manuscript lines 26-30
On peut supposer aussi (quitte à \unsure{restreindre|rétrécir} $X$) que \deleted{$n_x=\dim X_\eta$; \unsure{Posons, $n_x$} \unsure{et $n$}. Soit}
\[
\deleted{\unsure{n=\dim}}\qquad n=\dim X_\eta=n_x;
\]
Soit d'autre part
\[
r=\text{degtr. }k(x)/k.
\]
